\section{はじめに}
\label{sec:introduction}

テニス映像から選手の移動と打球の関係を分析するには、関節の運動、コート内での移動、ボールの位置を共通の時刻と座標で表す。
画像上の距離はカメラの向きと撮影距離に依存する。複数視点から奥行きを推定する際も、
白線の見切れ、選手の重なり、球の画面外移動で観測が途切れる。
クリップごとの修正は、復元を支えた観測の追跡と、多数のラリーへの同じ手順の適用を難しくする。

本報告は、同期済み固定カメラ映像を入力とする標準再構成を、幾何推定と実行システムの両面から記述する。
第\ref{sec:problem-and-output}章と第\ref{sec:reconstruction}章で、欠測を含む出力形式と再構成手順を示す。
第\ref{sec:court-detection}章から第\ref{sec:ball-detection}章では、合成教師、選手選別、球像と時間窓を比較し、
観測が失われる過程を調べる。第\ref{sec:execution-system}章と第\ref{sec:experiments}章で、実行システムと実映像での評価を報告する。
